\documentclass{article}

\usepackage[a-2u]{pdfx}
\usepackage[utf8]{inputenc}
\usepackage{array}
\usepackage[ddmmyyyy]{datetime}
\usepackage{geometry}
\usepackage{xcolor}
\usepackage{tikz}
\usetikzlibrary{shapes.geometric}

\renewcommand{\dateseparator}{.}
\renewcommand{\baselinestretch}{1.5}

\geometry{margin=2cm}

\def\ResearchProject{Research project}
\def\CompanyProject{Company project}
\def\SoftwareProject{Software project}

% METADATA SECTION START
\def\StudentsFullName{Bc. Martin Hubata}
\def\Obor{Computer Science -- Software and Data Engineering}
\def\ProjetTitle{Development foundation and processes for legacy transportation software modernization}
\def\ProjectType{\CompanyProject}
\def\SupervisorFullName{Mgr. Petr Škoda, Ph.D.}
\def\ConsultantsFullName{[Industry consultant from M-line software s.r.o. -- TBD]}
\def\ExpectedStart{1.6.2026}
\def\ExpectedEnd{1.3.2027}
% METADATA SECTION END

\title{\ProjetTitle}
\author{\StudentsFullName}

\begin{document}

\centerline{\Large \textbf{Proposal of company software project}}
\centerline{ \textbf{Department of Software Engineering}}
\centerline{ \textbf{Faculty of Mathematics and Physics, Charles University}}
\bigskip
{\noindent\begin{tabular*}{\textwidth}{ >{\raggedleft}m{4cm} l}
 {\bf Solvers:} & \StudentsFullName \\
 {\bf Study program:} & \Obor \\
 & \\
 {\bf Project title:} & \ProjetTitle \\
 {\bf Project type:} & \ProjectType \\
 & \\
 {\bf Supervisor:} & \SupervisorFullName \\
 {\bf Consultants:} & \ConsultantsFullName \\
 & \\
 {\bf Expected start:} & \ExpectedStart \\
 {\bf Expected end:} & \ExpectedEnd \\
\end{tabular*}}

\section{Introduction}

The author works for M-line software as a software engineer, primarily on the EDISON product line, integrated into the team's standard development practices. The ``academic team'' for the project is this M-line development team. M-line is a Czech software company developing the EDISON information system for the transportation and logistics domain (public transport, payroll, timetables, and many other related agendas) on top of an InterSystems Caché backend with a Windows desktop client. The company plans an incremental modernization of the legacy .NET Framework + WinForms client onto a more modern stack of .NET Core + WPF.

The core goal of the project is to establish modern processes and an architectural foundation that subsequent in-house work will reuse. A concrete software piece -- a modernized client shell -- serves as both a working product and as a demonstration of the processes and patterns being introduced.

\medskip
\noindent The expected outcomes of the project:

\begin{itemize}
  \item \textbf{Testing infrastructure.} Unit-testing setup on both .NET (xUnit, NSubstitute) and Caché (native \%UnitTest), establishing a test-first culture absent from the legacy stack. Concretely, all public APIs in the toolkit and shell ship with unit tests, test patterns are documented for reuse by future agenda implementations, and automated test execution is part of the development workflow.
  \item \textbf{Architectural foundation and reusable toolkit.} Front-end components and code abstractions for future agenda implementations, including a backward-compatibility layer that hosts legacy agendas.
  \item \textbf{Licensing data migration.} Moving licensing configuration from an XML file into the Caché backend, with the related internal tooling created or rewritten. Preparation for reuse across other M-line products is in scope; implementation in those products is not.
  \item \textbf{Both client and server-side update mechanism modernization.} Most likely a full recreation; the existing mechanism is too tangled to develop incrementally.
  \item \textbf{Modernized shell with mini-app subsystem.} A new shell preserving the responsibilities of the legacy one but designed for maintainability and extensibility, exposing a mini-app interface that agendas can opt into for summaries and reminders displayed directly in or by the shell. At least one demonstrative mini-app will be delivered.
  \item \textbf{Pilot deployment.} Rollout on developer workstations first, then parallel deployment to users, and eventual replacement of the old shell. Modernization of agendas themselves is out of scope. Concretely, success means all M-line developer workstations using the new shell as their default daily tool, at least one customer pilot site running the new shell in parallel with the legacy, and all (still relevant) legacy agendas remaining launchable through the new shell.
\end{itemize}

\subsection{Relation to other projects}

The project is being developed alongside several existing M-line projects with which it must coexist:

\begin{itemize}
  \item \textbf{Legacy EDISON} (VB/C\# .NET Framework WinForms) -- the production system currently running at customer sites. The modernized version will run against the same Caché backend, share the general idea of responsibilities and will incrementally take over selected agendas. No breaking changes to the legacy contract are permitted during the transition.
  \item \textbf{InterSystems Caché backend} -- the shared data store. It is used by most of the company's projects. Many of the project's architectural decisions will be influenced by the fact that our deployment process is much faster when deploying just a Caché update (and not the .NET). This is impossible to change, due to the non-negotiable compatibility with the old agendas. A concrete example of this is the licensing. However, many lower-level decisions, like those in the toolkit for developers, will also be influenced by this fact. 
  \item \textbf{TemplateDesigner} -- a WPF agenda concerning documents (printable timetables, invoices, etc.), already written by a colleague partially in the new stack. It serves as a second consumer validating the toolkit design (alongside the shell itself); its needs will influence which parts of the toolkit are prioritized.
\end{itemize}

\definecolor{c4person}{HTML}{08427B}
\definecolor{c4container}{HTML}{1168BD}
\definecolor{c4external}{HTML}{999999}

\begin{figure}[h]
\centering
\begin{tikzpicture}[
  scale=0.75, transform shape, font=\small,
  c4person/.style={
    rectangle, rounded corners=6pt, fill=c4person, draw=none,
    minimum width=3cm, minimum height=2cm,
    text=white, align=center, inner sep=4pt,
    append after command={
      \pgfextra{%
        \fill[c4person] ([yshift=0.3cm]\tikzlastnode.north) circle (0.4);
      }
    }
  },
  c4container/.style={
    rectangle, rounded corners=4pt, fill=c4container, draw=none,
    text=white, align=center,
    minimum width=3.6cm, minimum height=2cm, inner sep=4pt,
  },
  c4external/.style={
    rectangle, rounded corners=4pt, fill=c4external, draw=none,
    text=white, align=center,
    minimum width=3.6cm, minimum height=2cm, inner sep=4pt,
  },
  c4database/.style={
    cylinder, draw=c4external!50!black, shape border rotate=90, aspect=0.4,
    fill=c4external, text=white, align=center,
    minimum width=3.2cm, minimum height=2.8cm,
  },
  c4boundary/.style={draw, rounded corners=6pt, line width=0.8pt, fill=none},
  c4arrow/.style={-latex, thick, dashed, draw=gray!60, font=\scriptsize, align=center}
]
  \node[c4person] (user) at (-4, 10) {%
    \textbf{End user}\\{\footnotesize [Person]}\\{\scriptsize Daily operator}
  };

  % Outer EDISON 2 boundary (contains Core + TemplateDesigner)
  \node[c4boundary, minimum width=17cm, minimum height=9cm] (outer) at (1, 4) {};
  \node[anchor=north west, font=\itshape\small, align=left]
    at ([shift={(0.3cm, -0.3cm)}]outer.north west)
    {EDISON~2};

  % Inner EDISON 2 Core boundary
  \node[c4boundary, minimum width=12cm, minimum height=7cm] (core) at (-0.9, 4) {};
  \node[anchor=north west, font=\itshape\small, align=left]
    at ([shift={(0.3cm, -0.3cm)}]core.north west)
    {EDISON~2 Core};

  \node[c4container] (shell)   at (-4, 5.5) {%
    \textbf{Modernized Shell}\\{\footnotesize [.NET (Core) + WPF]}\\{\scriptsize New client app}
  };
  \node[c4container] (toolkit) at ( 2.5, 5.5) {%
    \textbf{Shared Toolkit}\\{\footnotesize [.NET (Core) + WPF library]}\\{\scriptsize Reusable UI and view-model features}
  };
  \node[c4container] (bridge)  at (-4, 2) {%
    \textbf{Legacy Bridge}\\{\footnotesize [.NET Framework WinForms]}\\{\scriptsize Hosts WinForms agendas}
  };

  % TemplateDesigner — peer of Core within EDISON 2 outer
  \node[c4container] (td) at (7.5, 2) {%
    \textbf{TemplateDesigner}\\{\footnotesize [.NET (Core) + WPF]}\\{\scriptsize Document design tool}
  };

  % Legacy EDISON boundary (contains Legacy Shell + Agendas)
  \node[c4boundary, minimum width=13cm, minimum height=4cm] (legacyboundary) at (-4, -3) {};
  \node[anchor=north west, font=\itshape\small, align=left]
    at ([shift={(0.3cm, -0.3cm)}]legacyboundary.north west)
    {Legacy EDISON};

  \node[c4external] (legacy) at (-8, -3) {%
    \textbf{Legacy Shell}\\{\footnotesize [.NET Framework WinForms]}\\{\scriptsize Old client app}
  };
  % Legacy Agendas — stack of WinForms .exe processes (3 rectangles offset upper-right)
  \node[rectangle, draw=c4external!60!black, line width=0.5pt, fill=c4external, rounded corners=4pt, minimum width=3.6cm, minimum height=2cm] at (0.3, -2.7) {};
  \node[rectangle, draw=c4external!60!black, line width=0.5pt, fill=c4external, rounded corners=4pt, minimum width=3.9cm, minimum height=2cm] at (0.15, -2.85) {};
  \node[c4external, draw=c4external!60!black, line width=0.5pt] (agendas) at (0, -3) {%
    \textbf{Legacy Agendas}\\{\footnotesize [.NET Framework WinForms]}\\{\scriptsize Business modules}
  };

  \node[c4database] (cache) at (0, -8) {%
    \textbf{InterSystems Cach\'e}\\{\footnotesize [Database]}\\{\scriptsize Shared data store}
  };

  \draw[c4arrow] (user.south) -- node[pos=0.4, right] {Uses\\\textit{[Desktop UI]}} (shell.north);
  \draw[c4arrow] (shell.east) -- node[midway, above] {Consumes\\\textit{[.NET reference]}} (toolkit.west);
  \draw[c4arrow] (toolkit.south west) -- node[midway, right] {Exposes\\agendas via} (bridge.north);
  \draw[c4arrow] (shell.south) -- node[midway, right] {Hosts\\agendas via} (bridge.north);
  \draw[c4arrow] (bridge.south) -- node[pos=0.50, above right=-2pt] {Spawns\\\textit{[cmdline params]}} (agendas.north west);
  \draw[c4arrow] (legacy.east) -- node[midway, above] {Spawns\\\textit{[.NET Remoting]}} (agendas.west);
  \draw[c4arrow] (td.north) -- node[pos=0.5, above right=-2pt] {Consumes\\\textit{[.NET reference]}} (toolkit.south east);
  \draw[c4arrow] (shell.south east) .. controls (4, 4) and (4, -5) ..
    node[pos=0.55, right] {Queries\\\textit{[HTTP]}} (cache.north east);
  \draw[c4arrow] (toolkit.south) -- node[pos=0.4, left=3pt] {Exposes agenda} (td.west);
  \draw[c4arrow] (legacy.south east) -- node[pos=0.5, left=2pt] {Queries\\\textit{[HTTP]}} (cache.west);
  \draw[c4arrow] (td.south) -- node[pos=0.5, right] {Queries\\\textit{[HTTP]}} (cache.east);
  \draw[c4arrow] (agendas.south) -- node[pos=0.3, right] {Queries\\\textit{[HTTP]}} (cache.north);
\end{tikzpicture}
\caption{C4 Container diagram for the EDISON product line. EDISON~2 Core (modernized shell, shared toolkit, legacy bridge) is the focus of this project; TemplateDesigner is a peer WPF agenda within EDISON~2 that also consumes the toolkit. The legacy bridge spawns existing WinForms agendas, which are also launched by the legacy shell. All systems share the InterSystems Cach\'e backend.}
\label{fig:c4-container}
\end{figure}

\section{Project description}

\textbf{End users and functionality.} EDISON installations are operated primarily by Czech transportation and logistics companies, primarily public-transport providers. Daily users are back-office staff in roles such as dispatching, payroll administration, ticketing, and accounting. The system covers the operational lifecycle of a transport company, with currently around thirty agendas corresponding to specific business processes (fleet operations, scheduling, payroll, ticketing, asset management, subsidies, accounting, reporting, etc.).

\medskip
\noindent \textbf{Stakeholders and value.} The primary stakeholder is M-line software, the company developing and maintaining EDISON; the resulting foundation will be reused by in-house developers to migrate further agendas after the project ends. The secondary stakeholders are the customer organizations operating EDISON installations, who will benefit from the overall improvements in the development process (speed, testing) and will receive a more modern user experience.

\medskip
\noindent The new foundation and toolkit will make (re)implementing new agendas easier in a number of ways:
\begin{itemize}
    \item Testability reduces manual regression checking and protects against breaks during refactoring.
    \item The architecture will (at least partially) enforce testability and responsibility separation. This separation is especially murky in parts of the old codebase, which makes modifications to it very difficult.
    \item Mini-apps let agendas surface key state and tasks directly in the shell.
    \item The new stack will be much more supportive of AI coding assistants, both by the nature of the technologies used (WinForms UI definitions are a notorious problem) and the mentioned testability and architecture benefits.
\end{itemize}

Other than the benefits for the new agendas, the previously mentioned points that generally fall under the umbrella of the shell will be improved and will be rid of a number of pain points (which slow down development):
\begin{itemize}
    \item Modernization of licensing will allow us to extend the system for new features. But even the old functionalities will be made more user-friendly, with the users being:
    \begin{itemize}
        \item Developers, when they need to add or modify the relationships between modules that they have developed.
        \item Licenser (company employee creating/modifying licenses), who will be able to use a more modern tool to express the constraints coming from agreements with end users more easily. The tool will be more in line with other tools the company has recently developed.
    \end{itemize}
    \item Update mechanism modernization will allow improvements, or at least prepare the systems for it, because the old one is simply too tangled to be modified.
    \item The old shell has a number of tangled issues with user (in)activity which are effectively impossible to solve, and negatively impact the systems of licensing, updating and the agendas themselves.
\end{itemize}

\medskip
\noindent \textbf{AI-augmented development.} The project is developed with extensive use of AI coding assistants throughout. Rather than treating AI tooling as a passive convenience, the project will explicitly document the patterns and guardrails that make it productive: project memory in the form of in-repo guidance files, verification pipelines before accepting generated changes, and a preference for explicit and strongly typed dependency injection over hidden global states. This will increase the quality of the generated code and improve the whole process, including the code review.

\medskip
\noindent \textbf{Delivery and defense format.} The academic deliverable consists of:
\begin{itemize}
  \item \textbf{Source code} of the modernized client shell, the supporting toolkit library, the demonstrative mini-app, the supporting test projects, and any internal tools modified or created during the project where the changes are substantial.
  \item \textbf{Developer documentation}: architectural overview, toolkit usage guide, and recommended development setup for future agenda implementations.
  \item \textbf{AI documentation}: patterns, guardrails, and observations collected during the project, delivered either as a standalone document or integrated with the developer documentation.
  \item \textbf{User documentation} for the new shell, written for end users and primarily highlighting new features and differences from the old shell.
\end{itemize}
These are delivered as a snapshot from M-line's GitHub repositories (not the full repositories themselves); the academic supervisor, opponent, and Committee receive access on request. The project is also presented as a live demonstration at the defense.

\medskip
\noindent \textbf{Company consent.} M-line software s.r.o. has approved the project scope and consents to its academic delivery and presentation as described above.

\end{document}
